\section{Увод}

Електронната търговия е сред най-динамично развиващите се области на съвременните информационни системи. За разлика от ранните уеб каталози, които изпълняват предимно представителни функции, днешните онлайн магазини представляват комплексни приложения, в които се преплитат процеси по идентификация на потребители, представяне на продукти, управление на наличности, обработка на поръчки, интеграция с външни услуги и устойчиво съхранение на бизнес данни. В този контекст разработването на уеб базирана система за електронен магазин е подходяща инженерна задача, тъй като позволява едновременно разглеждане на архитектурни, алгоритмични, организационни и интерфейсни аспекти на един завършен софтуерен продукт.

Конкретната предметна област на настоящата работа е насочена към продажба на компютърна и офис техника. Този избор не е случаен. От една страна, подобни продукти са масово търсени както от индивидуални потребители, така и от фирмени клиенти. От друга страна, те предполагат ясно структуриране по категории, проследяване на наличности, асоцииране с доставчици или фирми, и поддържане на сравнително богата каталожна информация. По тази причина предметната област е достатъчно богата, за да наложи реални архитектурни решения, но и достатъчно ограничена, за да може да бъде обхваната в рамките на дипломна разработка.

Разработеният проект е реализиран като Java уеб приложение върху Spring екосистемата. Този избор позволява използване на добре утвърдени индустриални практики при изграждане на многослойни информационни системи: разделение между контролери, бизнес логика, достъп до данни и визуален слой; декларативна конфигурация на част от инфраструктурните компоненти; валидация на входните данни; и стандартен процес за компилация и пакетиране с Maven. В резултат на това проектът може да бъде разглеждан не само като конкретен онлайн магазин, но и като учебен пример за систематично приложение на основни принципи от софтуерното инженерство.

\subsection{Актуалност на темата}

Темата е актуална по няколко причини. На първо място, дигитализацията на търговските процеси продължава да измества традиционните форми на продажба към уеб и мобилни платформи. Потребителите очакват постоянен достъп до каталожна информация, лесно сравнение между продукти, бърз процес на поръчка и възможност за проследимост на взаимодействието със системата. Това поставя високи изисквания пред разработчиците на уеб приложения както по отношение на функционалното покритие, така и по отношение на архитектурната подреденост на решението.

На второ място, електронният магазин е тип система, в която дори относително малък проект обхваща множество фундаментални теми: моделиране на домейн обекти, изграждане на уеб маршрути, работа с шаблони, персистентност, защита на потребителски данни, валидация на форми, управление на състояние и основни интеграции с външна инфраструктура. Това прави задачата особено ценна в образователен контекст, защото позволява синтез на знания от различни дисциплини в една практическа разработка.

На трето място, специализацията към компютърна и офис техника придава конкретика на приложението. Наличието на продуктови категории като видеокарти, процесори, слушалки, принтери, офис столове, скенери, рутери, мишки и клавиатури прави системата достатъчно разнообразна, за да демонстрира работа с категоризация, филтриране, наличие и потребителски избор. Добавената поддръжка за фирми, които се използват при въвеждане на продукти, допълнително обогатява модела на системата и я доближава до реален бизнес сценарий.

\subsection{Проблемна постановка}

При изграждането на онлайн магазин възникват няколко взаимосвързани проблема. Необходимо е системата да предлага интуитивен потребителски интерфейс за разглеждане и избор на продукти, но едновременно с това да поддържа и надеждна вътрешна организация на данните. Потребителите трябва да могат да се регистрират, да влизат в профилите си, да добавят продукти в количка и да финализират поръчки. В същото време приложението трябва да управлява информация за фирми, категории и наличности, като запазва консистентност между визуалното представяне и реалното състояние на данните в базата.

Особено предизвикателство представлява нуждата от ясно архитектурно разделение. Ако цялата логика бъде концентрирана в контролерите или шаблоните, приложението бързо става трудно за поддръжка и разширяване. Затова е необходимо да се изгради решение, при което представянето, бизнес логиката и достъпът до данни да бъдат отделени в самостоятелни слоеве. Допълнително усложнение възниква от необходимостта за обработка на потребителски вход, валидиране на данни, защита на паролите и интеграция с електронна поща за потвърждение на регистрацията.

В рамките на настоящия проект проблемът може да бъде формулиран по следния начин: да се разработи компактна, но функционално завършена уеб система за електронна търговия, която да демонстрира цялостен процес на работа с потребители, фирми, продукти, количка и история на поръчките, като при това използва ясна многослойна архитектура и реални технологии от Java екосистемата.

\subsection{Цел на дипломната работа}

Основната цел на дипломната работа е да се проектира и разработи уеб базирана система за електронна търговия с компютърна и офис техника, която да поддържа основните процеси по регистрация на потребители, управление на фирми и продукти, разглеждане на каталог, работа с количка, финализиране на поръчка и проследяване на история на покупки.

Постигането на тази цел предполага не само създаване на работещ код, но и обосноваване на използваните архитектурни и технологични решения. С други думи, системата следва да бъде едновременно реализирана и аналитично описана така, че документът да показва не само какво прави приложението, но и защо е избран именно този подход.

\subsection{Задачи на разработката}

За реализиране на поставената цел са формулирани следните основни задачи:

\begin{enumerate}[leftmargin=*]
    \item анализ на предметната област и формулиране на основните потребителски сценарии;
    \item избор на подходящ технологичен стек за създаване на сървърно рендирано Java уеб приложение;
    \item проектиране на многослойна архитектура, която разделя визуален слой, контролери, услуги и достъп до данни;
    \item моделиране на домейн обектите за потребители, фирми, продукти, количка и история на поръчки;
    \item реализация на регистрация, вход, сесийно управление и потвърждение на акаунт чрез електронна поща;
    \item реализация на механизми за добавяне и категоризация на продукти, включително асоциирането им с фирма;
    \item изграждане на каталог с визуализация на продукти, преглед на детайли и добавяне в количка;
    \item реализиране на процес по оформяне на поръчка и архивиране на историята на покупките;
    \item реализиране на базова локализация на интерфейса и примерен REST интерфейс за агрегирана информация;
    \item анализ на структурата на проекта, конфигурационните файлове и генерираните build артефакти;
    \item оценка на силните страни, ограниченията и възможностите за бъдещо развитие на системата.
\end{enumerate}

\subsection{Обект, предмет и обхват на изследването}

Обект на изследване е процесът по организиране и реализиране на електронна продажба на компютърна и офис техника чрез уеб приложение. Предмет на изследване са архитектурните, технологичните и програмните средства, чрез които този процес може да бъде моделиран и имплементиран в рамките на Java базирана информационна система.

Обхватът на разработката включва следните аспекти:

\begin{itemize}[leftmargin=*]
    \item back-end реализация на основните бизнес процеси в приложението;
    \item уеб интерфейс, реализиран чрез HTML шаблони и статични ресурси;
    \item персистентност на данни чрез MySQL и JPA модел;
    \item локална конфигурация за връзка с база данни и SMTP сървър;
    \item използване на \texttt{src} като директория за сорс код и \texttt{target} като директория за генерирани артефакти.
\end{itemize}

Извън обхвата на настоящата работа остават няколко характерни за производствените системи елемента: интеграция с реален платежен оператор, роля-базирана система за авторизация чрез Spring Security, контейнеризация с Docker, централизирано управление на тайни и разширена система от автоматизирани тестове. Тези аспекти са важни, но съзнателно са оставени като естествено поле за бъдещо развитие, тъй като основната цел тук е изграждането на функционална и добре структурирана учебно-приложна система.

\subsection{Методология на работа}

Подходът към разработката съчетава елементи от анализ, проектиране, имплементация и последваща оценка. Първо е разгледана предметната област и са извлечени основните сценарии, които системата следва да покрива. След това е избран технологичен стек, който е достатъчно зрял, добре документиран и подходящ за изграждане на многослойно Java уеб приложение. Следващата стъпка е моделиране на ключовите домейн обекти и дефиниране на връзките между тях.

При реализацията се следва практиката различните отговорности да бъдат разпределени в отделни пакети и слоеве. Контролерите обработват HTTP заявките и подготвят данните за визуализация. Услугите реализират бизнес логиката. Repository интерфейсите осъществяват достъп до базата данни. Шаблоните визуализират състоянието на системата пред крайния потребител. Допълнително се използват binding, service и view модели за ограничаване на зависимостите между отделните слоеве.

Оценката на решението в рамките на настоящия документ е предимно аналитична и сценарийна. Тя стъпва върху изследване на реалния код, маршрути, шаблони, конфигурация и данни, като се формулират както постигнатите резултати, така и ограниченията на текущата имплементация. Този подход е особено подходящ за дипломна документация на реален проект, защото позволява да се анализира не само външното поведение, но и вътрешната логика на системата.

\subsection{Практическа стойност на разработката}

Практическата стойност на проекта се изразява в няколко направления. На първо място, създадена е реално стартираща основа за онлайн магазин, която може да бъде използвана като демонстрационен проект, учебна база или отправна точка за по-нататъшно надграждане. На второ място, системата демонстрира как могат да бъдат комбинирани ключови Spring технологии в рамките на един последователен архитектурен модел. На трето място, документацията осигурява систематично описание на решенията в проекта, което улеснява поддръжката, разбирането и бъдещото развитие.

В образователен план проектът има допълнителна стойност, защото показва как абстрактни понятия като MVC, dependency injection, ORM, валидация и транзакционност се материализират в конкретни класове, маршрути, шаблони и конфигурации. Това превръща системата в удобен мост между теорията и практиката.

\subsection{Структура на изложението}

Изложението е организирано така, че да следва логиката на една завършена инженерна разработка. След резюмето настоящата глава въвежда темата, мотивацията и задачите. Втората глава позиционира проекта спрямо съществуващи решения и утвърдени подходи в областта на електронната търговия. Третата глава представя теоретичните основи, необходими за разбиране на архитектурния и технологичния избор. Четвъртата глава формулира изискванията към системата и описва основните сценарии на употреба.

Петата глава разглежда общата архитектура и организацията на кода. Шестата глава е посветена на домейн модела, релациите между обектите и трансфера на данни между слоевете. Седмата глава анализира back-end имплементацията по модули. Осмата глава разглежда представящия слой, интерфейсните решения и локализацията. Деветата глава описва структурата на проекта, конфигурацията на средата и build артефактите. Десетата глава представя потребителските сценарии и предвидените места за screenshots от реалната работа на приложението. Единадесетата глава е посветена на верификацията и оценката. Дванадесетата глава обобщава ограниченията и формулира насоки за бъдеща работа. Последната глава съдържа заключение с общ синтез на постигнатото.

Така структуриран, документът цели да покаже не само крайния резултат, но и последователния път от проблема и изискванията към архитектурата, реализацията и критичната оценка на системата.
